\begin{proposition}[7.3.18]
\label{IV.2.s7.3.18-fr}
Supongamos que $\boldsymbol P$ es una propiedad del primer \textup{(resp.} del
segundo) tipo definida a partir de una propiedad $\boldsymbol R$ que
satisface las condiciones $(\mathrm R'_{\mathrm I})$ y
$(\mathrm R_{\mathrm{II}})$, así como la condición siguiente:

$(\mathrm R_{\mathrm{III}})$ Para todo anillo local noetheriano \textup{completo}
$C$, si se pone $Z=\operatorname{Spec}(C)$, el conjunto
$U_{\mathbf{R}}(Z)$ \textup{(7.3.12)} es abierto en $Z$.

Entonces, si $A$ es un $\boldsymbol P$-anillo y se pone
$X=\operatorname{Spec}(A)$, el conjunto $U_{\mathbf{R}}(X)$ es
abierto en $X$.
\end{proposition}

\begin{proof}
En efecto, si $X'=\operatorname{Spec}(\widehat A)$ y $f:X'\to X$ es
el morfismo canónico, se tiene
$U_{\mathbf{R}}(X')=f^{-1}(U_{\mathbf{R}}(X))$ (7.3.12).
Como $f$ es fielmente plano y cuasicompacto, y por hipótesis
$U_{\mathbf{R}}(X')$ es abierto en $X'$, la conclusión se deduce
de (2.3.12).
\end{proof}
